
\section{Pair production by a photon in a magnetic field}

\SourcePage{428}{433}

The production of an electron-positron pair by a photon in a magnetic field and magnetobremsstrahlung radiation are two crossed channels of one and the same reaction. Therefore the amplitude $M_{fi}$ of the process of pair production is obtained from the amplitude of bremsstrahlung radiation simply by the substitutions
\begin{equation}
\varepsilon, \mathbf{p} \to -\varepsilon_+, -\mathbf{p}_+; \quad \varepsilon', \mathbf{p}' \to \varepsilon_-, -\mathbf{p}_-; \quad \omega, \mathbf{k} \to -\omega, -\mathbf{k}
\label{eq:91.1}
\end{equation}
(here $\varepsilon_-$, $\mathbf{p}_-$ and $\varepsilon_+$, $\mathbf{p}_+$ are the energies and momenta of the electron and positron in the pair; $\varepsilon$, $\mathbf{p}$ and $\varepsilon'$, $\mathbf{p}'$ are the initial and final energies and momenta of the electron in bremsstrahlung radiation)\footnote{This circumstance is, however, clear in advance: the axial vector of the arising polarisation can only be directed along the sole axial vector $\mathbf{H}$ figuring in the problem.}. In terms of angles and absolute values the transformation of momenta is
\begin{equation}
|\mathbf{p}| \to |\mathbf{p}_+|, \quad |\mathbf{p}'| \to |\mathbf{p}_-|, \quad \theta \to \pi - \theta_+, \quad \theta' \to \theta_-, \quad \varphi \to \varphi - \pi,
\label{eq:91.2}
\end{equation}
where $\theta_\pm$ are the angles between the momenta of the electron and positron and $\mathbf{k}$, and $\varphi$ is the angle between the planes $\mathbf{k}$, $\mathbf{p}_+$ and $\mathbf{k}$, $\mathbf{p}_-$. In the case of bremsstrahlung radiation the cross section is expressed

\SourcePage{429}{434}

through the amplitude by the formula\footnote{In this section we again set not only $c = 1$, but also $\hbar = 1$.}
\begin{equation}
d\sigma = |M_{fi}|^2\,\frac{1}{8|\mathbf{p}|\varepsilon'\omega}\,\delta(\varepsilon - \varepsilon' - \omega)\,\frac{d^3p'\,d^3k}{(2\pi)^5}
\label{eq:91.3}
\end{equation}
(see (64.25)); the $\delta$-function is removed by integration over $\varepsilon'$. Noting that
\[
d^3p' = |\mathbf{p}'|\varepsilon'\,d\varepsilon'\,do', \quad d^3k = \omega^2\,d\omega\,do_\mathbf{k},
\]
we must simply replace
\[
\delta(\varepsilon - \varepsilon' - \omega)\,d^3p'\,d^3k \to \omega^2|\mathbf{p}'|\varepsilon'\,do_\mathbf{k}\,do'\,d\omega.
\]

Then
\begin{equation}
d\sigma = |M_{fi}|^2\,\frac{\omega|\mathbf{p}'|}{8(2\pi)^5|\mathbf{p}|}\,do_\mathbf{k}\,do'\,d\omega.
\label{eq:91.4}
\end{equation}

In the case of pair production by a photon the cross section is expressed through the amplitude as
\[
d\sigma = |M_{fi}|^2\,\frac{1}{8\omega\varepsilon_+\varepsilon_-}\,\delta(\omega - \varepsilon_+ - \varepsilon_-)\,\frac{d^3p_+\,d^3p_-}{(2\pi)^5}
\]
or, after elimination of the $\delta$-function:
\begin{equation}
d\sigma = |M_{fi}|^2\,\frac{|\mathbf{p}_+||\mathbf{p}_-|}{8(2\pi)^5\omega}\,do_+\,do_-\,d\varepsilon_+.
\label{eq:91.5}
\end{equation}

Comparing with (91.4), we see that to obtain the cross section of pair production from the cross section of bremsstrahlung radiation one must carry out the substitution (91.1) and multiply the latter by
\begin{equation}
\frac{\mathbf{p}_+^2\,d\varepsilon_+}{\omega^2\,d\omega}
\label{eq:91.6}
\end{equation}
and replace $do'\,do_\mathbf{k}$ by $do_+\,do_-$.

In the ultrarelativistic case ($\omega \gg m$)\footnote{More precisely, one should have $\omega\sin\vartheta \gg m$, where $\vartheta$ is the angle between $\mathbf{k}$ and $\mathbf{H}$; when $\vartheta = 0$ pairs are not produced at all. Below we set $\vartheta = \pi/2$.} this can be done in the formulas obtained in the preceding section. It is assumed that both particles of the pair are ultrarelativistic; it is easy to verify that in such a case all the approximations used in \S\,90 remain valid.

In particular, the probability of pair production by an unpolarised photon, summed over the spin projections of the electron and positron and integrated over the emission directions of the electron, is obtained by performing the substitution (91.1) in the formula (90.22) (more precisely, in the expression for $dI/\omega$), where $d^3k = \omega^2\,d\omega\,do_\mathbf{n}$ is replaced by $d^3p_+$:
\begin{equation}
\begin{split}
dw &= \frac{e^2}{4\pi^2}\,\frac{d^3p_+}{\omega}\int_{-\infty}^{\infty}\left(\frac{m^2}{\varepsilon_+\varepsilon_-} - \frac{\varepsilon_+^2+\varepsilon_-^2}{4\varepsilon_-^2}\omega_{0+}^2\tau^2\right) \times {} \\
&\quad \times \exp\!\left\{i\frac{\omega\tau\varepsilon_+}{\varepsilon_-}\left(1 - \mathbf{nv}_+ + \frac{\tau^2}{24}\omega_{0+}^2\right)\right\}d\tau,
\end{split}
\label{eq:91.7}
\end{equation}
where $\omega_{0+} = |e|H/\varepsilon_+$; $\mathbf{n}$ is the unit vector in the direction of the photon momentum, lying in the plane perpendicular to the magnetic field. The integration is performed in the same way as in \S\,90, since (expression (91.7)) depends only on the angle between $\mathbf{n}$ and $\mathbf{v}_+$---one can integrate indifferently over $do_+$ or over $do_\mathbf{n}$. The result can therefore be obtained directly by analogy with (90.23):
\begin{equation}
dw = \frac{m^2 e^2\,d\varepsilon_+}{\sqrt{\pi}\;\omega^2}\left[\int_x^\infty \Phi(\xi)\,d\xi + \left(\frac{2}{x} - \varkappa x^{1/2}\right)\Phi'(x)\right],
\label{eq:91.8}
\end{equation}
where now
\begin{equation}
x = \left(\frac{m^3\omega}{|e|H\varepsilon_+\varepsilon_-}\right)^{\!2/3}, \quad \varkappa = \frac{|e|H\omega}{m^3}\left(= \frac{\hbar^2|e|H\omega}{m^3c^5}\right).
\label{eq:91.9}
\end{equation}

\SourcePage{430}{435}

The total probability of pair production per unit time is obtained by integrating (91.8) over $\varepsilon_+$ (and, in view of the obvious symmetry with respect to $\varepsilon_+$ and $\varepsilon_- = \omega - \varepsilon_+$, it suffices to integrate from 0 to $\omega/2$ and then double the result). Performing the change of the integration variable from $\varepsilon_+$ to $x$ and integrating the first term in (91.8) twice by parts, we obtain
\begin{equation}
w = -\frac{|e|^3H}{m\varkappa\sqrt{\pi}}\int_{(4/\varkappa)^{2/3}}^{\infty}\frac{2(x^{3/2}+1/\varkappa)}{x^{11/4}(x^{3/2}-4/\varkappa)^{1/2}}\,\Phi'(x)\,dx
\label{eq:91.10}
\end{equation}
(\emph{A.~I.~Nikishov, V.~I.~Ritus}, 1967).

In the limiting case of weak fields ($\varkappa \ll 1$) the essential values of $x$ in the integral (91.10) are those near the lower limit. Since these values are large, one can use the asymptotic expression for the Airy function
\[
\Phi(x) \approx \frac{1}{2x^{1/4}}\exp\!\left(-\frac{2}{3}x^{3/2}\right)
\]
(see III, \S\,b). Introducing the variable $y = x^{3/2} - 4/\varkappa$ and setting $y = 0$ everywhere possible, we obtain in the result
\begin{equation}
w = \frac{3^{3/2}|e|^3H}{2^{5/2}m}\exp\!\left(-\frac{8}{3\varkappa}\right), \quad \varkappa \ll 1.
\label{eq:91.11}
\end{equation}
The exponential decrease of the probability when $\varkappa \to 0$ corresponds to the impossibility of pair production in the classical limit.

\SourcePage{431}{436}

In the opposite limiting case of strong fields ($\varkappa \gg 1$) only the second term is essential in the integral (91.10), and it is determined by the region of values of $x$ in which $x^{3/2} \sim 1/\varkappa \ll 1$. In this region one can replace $\Phi'(x)$ by its value
\[
\Phi'(0) = -\frac{3^{1/6}\Gamma(2/3)}{2\pi^{1/2}}.
\]

Using the value of the integral
\[
\int_1^\infty y^{-\nu}(y-1)^{\mu-1}\,dy = \frac{\Gamma(\nu-\mu)\Gamma(\mu)}{\Gamma(\nu)},
\]
we obtain
\begin{equation}
w = \frac{3^{1/6}\cdot 5\Gamma^2(2/3)}{2^{4/3}\cdot 7\pi^{1/2}\Gamma(7/6)}\,\frac{|e|^3H}{m\varkappa^{1/3}} = 0.38\,\frac{|e|^3H}{m\varkappa^{1/3}}, \quad \varkappa \gg 1.
\label{eq:91.12}
\end{equation}
The function $mw(x)/|e|^3H$ has a maximum of 0.11 at $\varkappa \approx 11$.
